\subsection{Randomized Stochastic Stress-Testing}
\label{sec:stochastic_eval}

While the deterministic micro-benchmark isolates policy transitions under a single controlled disturbance, operational 5G and 6G multi-backhaul deployments encounter non-stationary wireless fading, stochastic packet queueing, and unpredictable failure onsets. To rigorously evaluate the statistical stability of the $2 \times 2$ policy matrix, we subject all four steering engines to an automated 400-trial stochastic emulation benchmark ($N = 100$ independent trials per policy across a 20.0\,s session timeline).

\subsubsection{Stochastic Disturbance Model}
Each trial instantiates an identical stochastic profile across the four policies to guarantee strict paired comparability:
\begin{enumerate}
    \item \textbf{Stochastic Step Arrival}: The onset of permanent structural degradation on Link B is drawn from a uniform distribution:
    \begin{equation}
        t_{\text{step}} \sim \mathcal{U}(5.0, 9.0)\,\text{s}.
    \end{equation}
    At $t \ge t_{\text{step}}$, the router permanently steps Link B one-way delay by $+40$\,ms (yielding RTT $\approx 95.5$\,ms), establishing a permanent breach of the SLA threshold $\tau = 65.0$\,ms.
    \item \textbf{Poisson Micro-Burst Process}: In the pre-step window $t \in [0.5, t_{\text{step}} - 0.2]$\,s, transient noise bursts arrive via a Poisson process with rate $\lambda = 0.5\,\text{s}^{-1}$ (mean inter-arrival time $\mathbb{E}[\Delta t] = 2.0$\,s). Each burst features a sub-tick duration $d_{\text{burst}} \sim \mathcal{U}(50, 80)$\,ms ($< T_{\text{poll}} = 100$\,ms) and delay excursion $A_{\text{burst}} \sim \mathcal{U}(28.0, 42.0)$\,ms, driving instantaneous RTT to $78\text{--}92$\,ms ($> \tau$). Because $d_{\text{burst}} < T_{\text{poll}}$, isolated bursts cannot straddle multiple polling ticks, isolating single-sample noise from persistent channel decay.
\end{enumerate}

\paragraph{Physical Grounding \& Standards Alignment}
The emulation parameters reflect operational 5G and satellite backhaul characteristics:
\begin{itemize}
    \item \textbf{Telemetry Protocol Realism}: Active 10\,Hz probing ($T_{\text{poll}} = 100$\,ms) mirrors standardized Bidirectional Forwarding Detection (BFD, RFC 5880), TWAMP (RFC 5357), and 3GPP Performance Measurement Function (PMF) profiles, consuming negligible control-plane overhead ($10.24$\,kbps, $<0.01\%$ of link capacity).
    \item \textbf{Transient Disturbance Realism}: The sub-tick micro-bursts ($50\text{--}80$\,ms, $+28\text{ to }+42$\,ms delay excursion) replicate physical-layer wireless transients, including cellular HARQ retransmissions (where up to 4 MAC retransmission attempts add $20\text{--}60$\,ms of transient latency), momentary base-station buffer queueing, and beam-phasing realignment in LEO/mmWave links.
    \item \textbf{Structural Failure Realism}: The permanent $+40$\,ms delay step replicates genuine backhaul decay, such as satellite rain fade triggering Adaptive Modulation and Coding (AMC) downshifts (e.g., from 64-QAM to QPSK) or upstream carrier protection rerouting over higher-latency transit paths.
    \item \textbf{3GPP Handover Alignment}: Policy 2's two-tick persistence window ($K=2$ at 10\,Hz, confirming for $200$\,ms) aligns directly with the 3GPP Time-To-Trigger (TTT) specification (TS 38.331: $100\text{--}320$\,ms) established to prevent handover ping-pong.
\end{itemize}

\begin{figure}[htbp]
\centering
\includegraphics[width=0.48\textwidth]{../plots/stochastic_tradeoff_compact_table.pdf}
\caption{Empirical trade-off between false handover ratio and SLA violation detection delay $T_{\text{detect}}$ across 400 stochastic evaluation trials ($N=100$ runs per policy). The embedded table provides exact $\text{Mean} \pm \text{Std}$ metrics. Policy 2 suppresses $88\%$ of noise-induced false handovers (dropping to $10.0\%$) while achieving a bounded reaction delay of $120.6$\,ms on genuine structural failure.}
\label{fig:stochastic_tradeoff}
\end{figure}

\begin{table}[htbp]
\caption{Randomized Stress-Testing: Empirical Stochastic Aggregates ($N = 100$ Runs per Policy)}
\label{tab:random_comparison}
\centering
\begin{tabular}{lcccc}
\toprule
\textbf{Evaluated} & \textbf{False Handover} & \textbf{Detection Delay} & \textbf{Handovers} & \textbf{Mean SLA Viol.} \\
\textbf{Policy} & \textbf{Ratio (\%)} & \textbf{$T_{\text{detect}}$ (ms)} & \textbf{per Run} & \textbf{Duration (ms)} \\
\midrule
Policy 1 (Strict Raw)    & $88.0 \pm 32.5\%$          & $4 \pm 15$           & $1.0 \pm 0.0$          & $85 \pm 36$ \\
Policy 2 (Raw + Hyst)    & \textbf{$10.0 \pm 30.0\%$} & $121 \pm 60$         & \textbf{$1.0 \pm 0.0$} & \textbf{$329 \pm 112$} \\
Policy 3 (Window)        & $35.0 \pm 47.7\%$          & $67 \pm 67$          & $1.0 \pm 0.0$          & $251 \pm 173$ \\
Policy 4 (Win + Hyst)    & \textbf{$28.0 \pm 44.9\%$} & $143 \pm 100$        & $1.0 \pm 0.0$          & $341 \pm 217$ \\
\bottomrule
\end{tabular}
\end{table}

\subsubsection{Empirical Performance Analysis}
As compiled in Table~\ref{tab:random_comparison} and visualized in Fig.~\ref{fig:stochastic_tradeoff}:

\paragraph{Vulnerability of Memoryless Thresholding (Policy 1)}
Under Poisson channel noise ($\lambda = 0.5\,\text{s}^{-1}$ across a pre-step window $\Delta t \approx 6.5$\,s), the probability of encountering at least one micro-burst is $\mathbb{P}(\ge 1\,\text{burst}) = 1 - e^{-\lambda \cdot \Delta t} \approx 96.1\%$. Because Policy~1 evaluates raw samples with zero temporal memory ($K=1$), the very first transient spike encountered forces an immediate, false-positive route switch, resulting in an $88.0 \pm 32.5\%$ false handover ratio ($T_{\text{detect}} = 3.7 \pm 15.1$\,ms). This empirical collapse confirms that memoryless instantaneous steering is fundamentally unsuited for stochastic wireless channels.

\paragraph{Noise Suppression via Temporal Hysteresis (Policy 2)}
Enforcing a two-tick persistence constraint ($K=2$) reduces the false handover ratio from $88.0\%$ down to $10.0 \pm 30.0\%$, achieving an approximate $9\times$ suppression of noise-induced route flapping. Policy 2 only misfires when two independent Poisson micro-bursts arrive in immediate succession within $100$\,ms. When structural link degradation occurs at $t_{\text{step}}$, Policy 2 confirms the failure across two consecutive probes and executes route replacement within $120.6 \pm 60.2$\,ms. This stochastic result directly validates the deterministic Hero Trace benchmark ($T_{\text{detect}} = 121.8$\,ms), establishing that Policy 2 restores bounded transport latency within $\approx 1.2$ polling intervals with negligible memory overhead ($O(1)$ integer counter).

\paragraph{Moving Average Window Distortions (Policies 3 \& 4)}
Window-based smoothing (Policy 3, $N=3$) achieves intermediate noise suppression ($35.0 \pm 47.7\%$ false handovers) with $66.5 \pm 67.1$\,ms detection lag. However, moving averages suffer from FIFO buffer pollution: an isolated high-amplitude outlier remains resident in the queue for $3 \times T_{\text{poll}} = 300$\,ms, elevating the moving baseline. If mild channel jitter arrives while the buffer is contaminated, $\mu_{\mathcal{W}}$ breaches $\tau$. Combining window smoothing with hysteresis (Policy 4) lowers false alarms to $28.0 \pm 44.9\%$, but compounds detection lag to $142.9 \pm 100.3$\,ms, increasing user exposure to degraded channel conditions (cumulative SLA violation duration $341.0 \pm 217.3$\,ms).
